%% !TEX root = manuscript.tex

\chapter*{Conclusion}

Cette deuxième année d'alternance a été dominée par des enjeux d'ingénierie des données et
d'architecture système, avant de revenir au second semestre à la modélisation statistique.

Le premier semestre a été celui de la reconstruction du socle. La migration de Grimmo vers
l'ERP SPO ne consistait pas à transférer des données, mais à rétro-concevoir un modèle dont
les règles n'étaient pas documentées, à résoudre des problèmes d'intégrité référentielle avec
la gestion des identifiants dynamiques et du versionnement, et à adapter la technique aux
impératifs métier --- ce qu'illustre l'adoption des «~coquilles vides~» pour sécuriser le
démarrage de la facturation. En parallèle, j'ai maintenu le service sur le reporting, avec la
mise en œuvre d'une règle juridique complexe malgré les contraintes de nos outils, et
participé au déploiement du SSO.

Le second semestre a prolongé ce travail, puis l'a valorisé. La reprise des tiers a achevé de
constituer le socle commercial du nouvel ERP, en créant 1~434 clients acquéreurs sans doublon
détecté par les trois clés d'unicité retenues. Le Copilote Financier a ensuite exploité ce
socle pour outiller la direction financière sur trois questions de gestion, de bout en bout,
de la compréhension des exports jusqu'à une application que les contrôleurs de gestion ouvrent
eux-mêmes.

Ce parcours aura démontré la thèse annoncée en introduction, et il l'aura fait deux fois. Une
première fois en creux, lorsque la charge de la migration a repoussé le projet de modélisation
au second semestre : modéliser des phénomènes financiers sur une base fragmentée entre
l'ancien et le nouveau système aurait donné des résultats biaisés et non pérennes. Une seconde
fois en plein, lorsque le projet de Data Science a donné des résultats solides sur le rythme
d'écoulement, où les données avaient une profondeur temporelle, et modestes sur le risque de
marge, où elles n'étaient qu'une photographie à date. La structure de la donnée commande ce
que la modélisation peut en tirer : c'est le même énoncé, obtenu par deux chemins
indépendants.

Il reste beaucoup à faire. La migration n'est pas achevée, l'historique budgétaire reste à
reprendre, l'intégration à l'infrastructure Nexity s'annonce, et le Copilote Financier ne
deviendra un véritable outil de suivi que lorsque les exports auront été historisés sur une
durée suffisante. Ce que cette année aura établi n'en dépend pas. Un projet de données ne se
joue pas au moment du choix de l'algorithme mais bien avant, lorsqu'on décide ce qu'est un
client, ce qu'est une marge, ou ce que contient réellement un fichier dont le nom annonce
autre chose. Structurer une donnée et la modéliser ne sont pas deux métiers successifs, mais
deux moments du même travail, et c'est le premier qui décide de ce que le second pourra en
tirer.
